% Architecture Comparison Table
% Generated by Q60-PAPER-ARTIFACTS-EXPORT

\begin{table}[htbp]
\centering
\caption{Architecture comparison across the three binary classification pipelines.
  All three share a frozen MobileNetV3-Large backbone; only the head/circuit differs.}
\label{tab:architecture}
\begin{tabular}{llllll}
\toprule
\textbf{Model} & \textbf{Backbone} & \textbf{Head / Circuit} & \textbf{Simulator} & \textbf{Layers / Depth} & \textbf{Activation} \\
\midrule
Classical CNN (Q56)
  & MobileNetV3-Large (frozen) + Linear(960$\to$128, frozen proj.)
  & 3-layer MLP
  & ---
  & FC(128$\to$16), FC(16$\to$8), FC(8$\to$2)
  & ReLU + Softmax \\
DV-QNN (Q57)
  & MobileNetV3-Large (frozen)
  & 4-qubit VQC
  & qcore DV (discrete)
  & depth$=1$, $\alpha{=}0.1$
  & Pauli-Z measurement \\
CV-QNN (Q58)
  & MobileNetV3-Large (frozen)
  & 2-mode CV circuit
  & qcore Gaussian (CV)
  & depth$=1$, squeeze cap$=1.5$
  & Homodyne measurement \\
\bottomrule
\end{tabular}
\begin{tablenotes}
\small
\item All backbones initialised from ImageNet-pretrained weights; backbone parameters are frozen during training.
\item DV-QNN uses a custom \texttt{medical\_ansatz}; CV-QNN uses \texttt{GaussianBackend} + \texttt{GVA} ansatz.
\item MLP head is the Q34A reference head used across all Phase 6b benchmarks.
\end{tablenotes}
\end{table}
